\chapter*{Abstract}
\label{cha:abtract}
\addcontentsline{toc}{chapter}{Abstract}

% research problem and objectives, literature (test, analyze, evaluate), present
% tense

% methods, past tense

% key results and arguments, present or past simple

% conclusions, implications, present simple

Linux supports a vast number of powerful tracing capabilities. Over the history
of the operating system, multiple frameworks and event sources have been
introduced for different use cases. Due to the size and complexity of the Linux
kernel, its tracing capabilities are often not well understood by developers and
are frequently associated with "black magic". An analysis of the Linux kernel
codebase was conducted, gathering information through source code, git commits,
mailing lists, technical articles, and conference talks. This work presents a
description of the tracing infrastructure provided by the Linux kernel version
7.2, focusing on traditional frameworks like Ftrace and Perf, event sources like
tracepoints and kprobes, and modern developments through eBPF. Finally, a real
use case for eBPF in containerized environments is studied, presenting Kivebpf
as an application leveraging modern Linux tracing infrastructure.
